\documentclass{article}
\usepackage[T1]{fontenc}
\usepackage{lmodern}
\usepackage{graphicx}
\usepackage{amsmath,amssymb}
\usepackage[a4paper,margin=2cm]{geometry}
\usepackage[absolute,overlay]{textpos}
\setlength{\TPHorizModule}{1mm}
\setlength{\TPVertModule}{1mm}
\usepackage[indonesian]{babel}
\usepackage{hyperref}
\setlength{\emergencystretch}{2em}
% unit: Halaman 22

\begin{document}

% === Halaman 22 (python/native) ===

22

• XLEX dipetakan menjadi th*t. 

• Kata yang cocok untuk th*t. adalah that. 

• Jadi kita memperoleh: E → a

• Hasil iterasi pertama:

 heVeTCSWPeYVaWHaVSReQMthaYVaOeaWHRtatePFaMVaWHKVSTYhtZe

theKeetPeJVSZaYPaRRGaReMWQhMGhMtQaReWGPSReHMtQaRaKeaTtM

JTPRGaVaKaeTRaWHatthattMZeTWAWSQWtSWatTVaPMRtRSJGSTVRea

YVeatCVMUeMWaRGMeWtMJMGCSMWtSJOMeQtheVeQeVetQSVSTWHKPaG

ARCStRWeaVSWeeBtVeZMtFSJtheKaGAaWHaPSWYSWeWeaVtheStheVt

heRGaPeRQeVeeBGeeHMWYPFhaVHaWHYPSRRFQMthaPPtheaCCeaVaWG

eSJKTVWMRheHYSPHtheQeMYhtSJtheMWReGtQaROeVFVeZaVAaKPeaW

HtaAMWYaPPthMWYRMWtSGSWRMHeVatMSWMGSTPHhaVHPFKPaZeNTCMt

eVJSVhMRSCMWMSWVeRCeGtMWYMt

\end{document}
